%% FeatureLiftBench: FSE-oriented condensed manuscript, revised 2026-09-17.
%% Scope: Python-150 benchmark; 150 tasks x 6 configurations (900 outcomes).
%% Current inputs and model configurations: paper_sources.json.
\documentclass[acmsmall,screen,review,anonymous]{acmart}

\usepackage{amsmath}
\usepackage{booktabs}
\usepackage{tabularx}
\usepackage{graphicx}
\usepackage{xspace}
\usepackage{microtype}
\usepackage{placeins}

\newcommand{\flb}{\textsc{FeatureLiftBench}\xspace}
\newcommand{\rres}{\textsc{RRES}\xspace}
\graphicspath{{figures/}}
\setlength{\emergencystretch}{2em}
% Keep floats near their discussion while allowing text above/below them.
\renewcommand{\topfraction}{0.85}
\renewcommand{\bottomfraction}{0.65}
\renewcommand{\textfraction}{0.15}
\renewcommand{\floatpagefraction}{0.75}
\setlength{\textfloatsep}{14pt plus 3pt minus 2pt}
\setlength{\floatsep}{12pt plus 2pt minus 2pt}
\setlength{\intextsep}{12pt plus 2pt minus 2pt}
\citestyle{acmauthoryear}

\begin{document}

\title[FeatureLiftBench: Behavior-Preserving Feature Lifting]{FeatureLiftBench: Benchmarking Behavior-Preserving Feature Lifting from Real Repositories}
\author{Anonymous Author(s)}

\begin{abstract}
Reusing software across project boundaries requires preserving behavior after
its original context is removed. We introduce \flb, a benchmark for
\emph{behavior-preserving feature lifting}: given an intact source repository
and a public behavioral contract, a coding agent constructs an independent
package evaluated without runtime access to the repository. The benchmark
contains 150 Python tasks from 126 repositories, validated through automated
checks, multi-author semantic review, and source-free reference replay.
We evaluate six model--harness configurations. A paired ablation on 40 tasks shows
that repository evidence improves Luna's and Qwen's success by 35.0 and 27.5
percentage points. 79.5\% of Primary/Extended-first failures contain confirmed
reads of entrypoint-associated source-file content. Qualitative analysis of a
purposively diverse 40-case sample identifies recurring but limited patterns involving
changed ambient assumptions, incompatible supporting representations and
pipelines, and destination-specific obligations requiring adaptation.
Three paired implementation cases illustrate these patterns through platform,
diagnostic, and resource relationships. These observations motivate verifying public
behavioral obligations together with the relationships that support them,
while distinguishing preserved upstream behavior from requirements introduced
by adaptation. \flb provides an executable setting for studying how repository
evidence is translated into behavior that survives a new software boundary.
\end{abstract}

\begin{CCSXML}
<ccs2012>
 <concept><concept_id>10011007.10011074.10011099.10011102</concept_id>
 <concept_desc>Software and its engineering~Software testing and debugging</concept_desc>
 <concept_significance>500</concept_significance></concept>
 <concept><concept_id>10011007.10011074.10011099</concept_id>
 <concept_desc>Software and its engineering~Software verification and validation</concept_desc>
 <concept_significance>300</concept_significance></concept>
</ccs2012>
\end{CCSXML}
\ccsdesc[500]{Software and its engineering~Software testing and debugging}
\ccsdesc[300]{Software and its engineering~Software verification and validation}
\keywords{coding agents, software reuse, feature extraction, behavioral preservation, benchmark}
\maketitle

\section{Introduction}
\label{sec:introduction}

Software reuse involves selecting, adapting, and integrating existing
artifacts~\cite{krueger1992reuse}. Moving a capability across project boundaries
is one such reuse scenario: a developer may need a parser from a larger tool, configuration
handling from a framework, or a signal registry for a new library. Although the
source implementation is available, the target capability may extend beyond any
single file or function: its behavior can depend on supporting modules, shared
state, framework mechanisms, configuration, and packaged resources. Reuse
therefore requires more than locating code. The relevant behavior must be
reconstructed so that it continues to hold after the original project context is
removed. This raises a software engineering question: \emph{which behavioral
obligations remain unmet when an existing capability is reconstructed across a
new package boundary, even with its original implementation available?} We
study this question through coding agents that construct independent packages
from intact repositories and explicit behavioral contracts.

We formulate this task as \emph{behavior-preserving feature lifting}, or
\emph{feature lifting} for short. Given a version-pinned source repository and a
public behavioral contract, an agent constructs a new package that is evaluated
without runtime access to the source repository. The agent may extract, adapt, or
reimplement code; correctness is determined by the observable behavior required
at the new package boundary. Feature lifting therefore couples repository
understanding with software reuse: the implementation is available as evidence,
but the resulting capability must survive outside the environment that originally
supported it. Figure~\ref{fig:pipeline} illustrates this setting.

The central reuse problem is preserving behavioral obligations whose
implementation may depend on relationships beyond the apparent entrypoint.
For adapted tasks, these obligations can also require an upstream capability to
operate under a changed interface or context. We therefore distinguish the
behavior required at the destination boundary from the particular implementation
and assumptions found in the source repository.

\begin{figure}[htb]
  \centering
  \includegraphics[width=\linewidth]{figures/fig1.png}
  \caption{Motivation and task overview of \flb. An agent receives the intact
  source repository and a public behavioral contract, constructs a new package,
  and submits only that package to source-free evaluation. The repository
  structure and highlighted locations are illustrative; agents receive no
  benchmark-provided source-location hints.}
  \Description{Three panels show a source capability supported by multiple
  repository components, reconstruction into a new package, and source-free
  evaluation of the submitted package.}
  \label{fig:pipeline}
\end{figure}

This setting differs from two common forms of repository-level evaluation.
SWE-bench and SWE-Bench Pro ask agents to modify an existing project in response
to an issue~\cite{jimenez2024swebench,deng2025swebenchpro}, while FeatBench and
FeatureBench study feature implementation in an existing or partially removed
codebase~\cite{chen2026featbench,zhou2026featurebench}. In feature lifting, by
contrast, the complete implementation of the target capability is intentionally
available as evidence, but the deliverable is a new independent package whose
behavior is evaluated without the donor repository. This also differs from
classical software transplantation, which extracts donor functionality for
integration into a supplied host program~\cite{barr2015transplantation}. Our
focus is the general coding-agent setting in which the destination boundary is a
standalone package and the extraction strategy is left open.

To study this problem, we introduce \flb, a benchmark of 150 Python tasks drawn
from 126 repositories and 132 pinned source snapshots. Each task provides the
complete source repository and a public behavioral contract, while withholding
benchmark tests, reference implementations, and source-location hints. Submitted
packages are evaluated through loading, behavioral, and source-isolation checks
without the source repository. Benchmark construction combines explicit
contracts and protected tests with automated consistency checks, AI-assisted
screening and author manual review of all 150 retained tasks, and repeated
execution of feasible reference implementations.

We use six model--harness configurations on the same 150 tasks to examine
behavioral preservation, the role of repository evidence, and the resulting
independent implementations. Functional pass rates range from 24.0\% to 76.7\%.
For the two strongest configurations, DeepSeek V4 Pro and Flash, all
150 tasks receive usable submissions, yet 76 of their 77 combined failures first
occur at Primary or Extended evaluation gates rather than during delivery or loading. A paired
source-evidence ablation further shows that intact repository evidence materially
improves reconstruction: GPT-5.6 Luna gains 35.0 percentage points and Qwen gains
27.5 points over Contract Only. However, source access is not sufficient by
itself. Among 303 Primary/Extended-first failures in the main comparison, 241 (79.5\%)
contain confirmed reads of entrypoint-associated source-file content. These observations motivate examining specific obligations and
the software relationships that support them. Paired artifact inspections
illustrate platform-dependent path construction, shared diagnostic
representations, and adaptation of resource traversal to a destination policy.
Secondary analyses characterize continued execution after an observed passing
checkpoint and the footprints of successful independent packages.

This work makes three contributions:
\begin{itemize}
  \item \textbf{A task formulation for behavior-preserving feature lifting.}
  We define a repository-level reuse setting in which intact implementation
  evidence is available during construction, but the requested capability must
  execute across a new, source-free package boundary.

  \item \textbf{FeatureLiftBench, an executable benchmark of real repository
  reuse.} We construct and validate 150 tasks from 126 Python repositories,
  combining bounded behavioral contracts, protected tests, source-free
  evaluation, multi-author semantic review, and repeated reference execution.

  \item \textbf{An empirical study of agent-mediated, behavior-preserving
  software reuse.} Across six configurations, we examine outcomes across reuse
  structures, the contribution of repository evidence through paired ablation,
  and obligation--relationship comparisons in cases with confirmed reads of
  entrypoint-associated source-file content. We additionally characterize
  execution after an observed passing checkpoint and the implementation
  footprints of successful independent packages.
\end{itemize}

\section{The FeatureLiftBench Benchmark}
\label{sec:benchmark}

\flb comprises 150 executable feature-lifting tasks. Each task connects an
existing capability in a complete source repository to a public contract for a
new package. The agent receives implementation evidence and the contract; the
submitted package is assessed by protected behavioral tests in a source-free
environment. A reference implementation provides evidence that the task is
feasible. We first define the task and its visibility boundary, then describe
construction, validation, and benchmark composition.

\subsection{Task Formulation and Running Example}

A \flb task consists of a pinned source repository $R$, a public behavioral
contract $S$, an initial workspace $W$, and a protected evaluator $J$. A model
backend $M$ acting through an agent harness $H$ produces a submitted artifact
$A$:
\begin{equation}
  A = \operatorname{Run}(M,H,R,S,W).
\end{equation}

We use a behavioral contract to make interface obligations
explicit~\cite{meyer1992contract}. It specifies the required API surface and observable behavior,
including relevant defaults, exceptions, state changes, ordering constraints,
and explicit exclusions. The agent may inspect the complete pinned repository,
including upstream tests and documentation, but must construct the result in a
designated submission directory. Benchmark tests, reference implementations,
test-to-contract mappings, and benchmark-provided source-location hints remain
hidden.

The required output is an independently executable package under the benchmark
environment. The package may extract, adapt, or reimplement upstream code, but
it must execute without runtime access to the source repository. The contract,
rather than any particular source slice, defines the preservation boundary.
Consequently, multiple structurally different implementations may satisfy the
same task.

Functional success is defined as
\begin{equation}
  F(A)=B(A)\land T_{\mathrm{primary}}(A)
       \land T_{\mathrm{extended}}(A)\land I(A),
  \label{eq:functional}
\end{equation}
where $B$ checks that the submitted package can be loaded under the frozen
evaluation environment, $T_{\mathrm{primary}}$ and
$T_{\mathrm{extended}}$ exercise protected behavioral obligations, and $I$
checks source independence. All four conditions must hold.

\paragraph{Running example: a signal registry.}
The Blinker task asks the agent to reconstruct \texttt{Signal} and
\texttt{Namespace} under a new package boundary. Its contract requires behavior
such as sender-specific dispatch, weak-receiver cleanup, stable namespace
identity, and scoped connection management, while explicitly excluding
asynchronous dispatch and the upstream global named-signal singleton. Thus, an
implementation whose central \texttt{connect} and \texttt{send} paths work can
still fail if supporting state or lifetime behavior is lost. The agent may reuse
or reimplement the relevant upstream logic, but the resulting package must
satisfy the contract without importing or reading the original \texttt{blinker}
source at runtime.

\subsection{Task Construction}
\label{sec:construction}

Construction turns an existing repository capability into a bounded,
source-backed lifting task with a public behavioral contract, protected tests,
and a feasible source-independent reference implementation.
Figure~\ref{fig:construction} summarizes the construction and validation
workflow.

\begin{figure}[htb]
  \centering
  \includegraphics[width=\linewidth]{figures/fig2.png}
  \caption{Construction and validation of \flb. Construction turns a pinned
  source capability into a public contract, protected tests, and a feasible
  source-independent reference package. Validation combines automated
  consistency checks, multi-author semantic review supported by structured
  AI-assisted checks, and three successful source-free reference replays per task
  (450/450 executions).}
  \Description{The upper band shows task construction from capability selection
  and source pinning through contract definition and protected evaluation assets.
  The lower band shows structural checks, semantic review, and repeated reference
  execution before tasks enter the frozen benchmark.}
  \label{fig:construction}
\end{figure}

\paragraph{Select a capability and pin its source.}
We select capabilities from Python libraries, developer tools, and framework
components when their observable behavior can be stated explicitly and executed
without live external services. Each task is tied to a pinned source snapshot.
Agents receive the complete tracked source tree, including upstream tests and
documentation, rather than a feature-specific slice. Identifying the relevant
implementation and its supporting dependencies therefore remains part of the
task.
Candidate tasks that could not be specified, evaluated, or validated under this
protocol were excluded before the benchmark was frozen.

\paragraph{Define the behavioral contract.}
For each capability, we specify the exported APIs and the observable behavior
that must be preserved, including relevant signatures, defaults, exceptions,
state changes, and explicit exclusions. The contract may also declare an
adaptation when the requested package interface intentionally differs from the
upstream capability. The resulting task statement describes what the package
must do without prescribing where the implementation resides or how it should
be reconstructed; benchmark-provided source-location hints are withheld.

\paragraph{Construct protected tests and a feasible reference.}
Primary and Extended test groups operationalize the same public contract:
Primary covers central obligations, while Extended exercises additional cases
and combinations without introducing new requirements. Both groups are hidden
from the agent. Test-to-contract mappings are maintained during construction to
check that evaluator expectations remain within the declared boundary.

For each task, we also retain a reference package that passes the same
source-free functional evaluation as agent submissions. References may be
constructed by extracting and adapting upstream code or by task-specific
implementation. They establish that the task is executable under the benchmark
environment and provide a reference point for artifact-size measurements;
they are not necessarily unique or minimal implementations.

\paragraph{Freeze the benchmark.}
After validation, task specifications, pinned source identities, protected
tests, reference packages, dependency policies, and evaluator settings are
versioned together. The frozen benchmark contains 150 tasks from 126
repositories and 132 source snapshots. The same 150-task membership is used for
every configuration in the main comparison.

\subsection{Benchmark Validation}
\label{sec:validation}

Validation must address whether evaluator judgments capture the intended
behavior, an instance of the test-oracle problem~\cite{barr2015oracle}.
We validate the frozen benchmark at three complementary levels: structural
consistency, semantic alignment, and executable feasibility.

\paragraph{Structural consistency.}
Automated checks verify required task assets, source mappings, dependency
declarations, generated task statements, and test-to-contract mappings. All 150
retained task packages pass these checks and resolve to 132 verified source
snapshots. These checks target inconsistencies such as stale specifications,
missing API declarations, and invalid test mappings.

\paragraph{Semantic alignment.}
AI-assisted checks were used across all task packages as structured review
support rather than as semantic judges. All 150 retained tasks were subsequently
manually reviewed by the authors for consistency among the public contracts,
protected tests, pinned-source evidence, explicit exclusions, and taxonomy
labels. Final retention and correction decisions were made by the authors.

\paragraph{Executable feasibility.}
Each retained reference package is replayed three times under the same
source-free functional evaluation used for submissions. All 450 executions
(150 tasks $\times$ 3 runs) pass. Reference replay therefore establishes that
every frozen task admits at least one executable solution under the benchmark
environment; it does not by itself establish semantic completeness, which is
addressed separately by contract and test review.



\subsection{Dataset Composition}
\label{sec:suite}

We characterize the 150 tasks along three complementary dimensions:
\emph{functional family} describes what capability is lifted, \emph{lift type}
describes how the requested capability relates to its upstream implementation,
and \emph{entanglement} describes which forms of project dependency must be
handled when crossing the new package boundary. These annotations describe
benchmark composition and do not affect scoring.

\paragraph{Functional coverage.}
The benchmark spans ten functional families across 126 repositories and 132
pinned source snapshots. The largest families are parsing/tokenizing/decoding
(31 tasks), registry/plugin dispatch (22), serialization/formatting/rendering
(19), configuration resolution/discovery (18), and
validation/normalization/construction (15). The remaining tasks cover workflow
and session orchestration, resource and metadata loading, algorithms and data
structures, protocol state transitions, and cache/retry policies.

\paragraph{Lift types.}
Each task receives one of three labels according to the relation between the
requested capability and its upstream implementation. \emph{Direct} tasks
preserve an identifiable upstream capability while extracting it from its
original project context. \emph{Adapted} tasks retain an identifiable capability
but require a declared change to its interface or observable behavior.
\emph{Composite} tasks combine multiple upstream capabilities or reconstruct a
workflow without a single corresponding upstream API. The benchmark contains 56
Direct, 76 Adapted, and 18 Composite tasks. These labels characterize the
requested transformation rather than the implementation strategy used by an
agent.

\paragraph{Entanglement mechanisms.}
Tasks may additionally carry multiple labels describing the dependencies that
must be accounted for during lifting. We group the underlying annotations into
four mechanisms: \emph{code dependencies}, covering internal helpers, transitive
imports, and third-party interfaces; \emph{data and state}, covering shared
models, parser state, and registries; \emph{framework mechanisms}, covering
lifecycle, dispatch, and dynamic discovery; and \emph{environment and
resources}, covering configuration, path rules, and packaged resources. These
mechanisms are non-exclusive and describe where a capability is coupled to its
host project rather than independent difficulty factors.

Figure~\ref{fig:coverage} summarizes the resulting composition. Panel~(a) shows
the ten functional families partitioned by lift type, while panel~(b) shows the
prevalence of the four entanglement mechanisms overall and within each lift
type.

\begin{figure}[htb]
  \centering
  \begin{minipage}[c]{0.57\linewidth}
    \centering
    \includegraphics[width=\linewidth]{figures/fig3a_families.pdf}
  \end{minipage}
  \hfill
  \begin{minipage}[c]{0.41\linewidth}
    \centering
    \includegraphics[width=\linewidth]{figures/fig3b_entanglement.pdf}
  \end{minipage}
  \caption{Composition of the 150-task benchmark. (a) Ten functional families
  partitioned by lift type: Direct (56), Adapted (76), and Composite (18).
  (b) Coverage of four non-exclusive entanglement mechanisms overall and within
  each lift type; cells report percentages and task counts. These annotations
  characterize benchmark composition rather than predefined difficulty.}
  \Description{Panel (a) shows ten functional families as stacked horizontal
  bars partitioned by the three lift types. Panel (b) shows a heatmap of code
  dependencies, data and state, framework mechanisms, and environment or
  resource entanglement, overall and within each lift type.}
  \label{fig:coverage}
\end{figure}

\section{Evaluation Protocol and Experimental Setup}
\label{sec:protocol}

We evaluate six model--harness configurations on the same 150 tasks under a
shared information and evaluation policy. This section specifies the evaluated
systems, functional and artifact metrics, statistical protocol, and analyses of
repository evidence, source exposure, and execution effort.

\subsection{Information Boundary and Agent Configurations}

We evaluate DeepSeek V4 Pro (\texttt{deepseek-v4-\allowbreak pro-0813}), DeepSeek V4 Flash
(\texttt{deepseek-v4-\allowbreak flash-0731}), GPT-5.6 Luna, GLM-5.3-Flash,
Qwen3.6-35B-A3B, and GPT-OSS 120B, abbreviated as Pro, Flash, Luna, GLM,
Qwen, and OSS. All configurations use OpenHands CLI 1.16.0 with native tool
calling~\cite{wang2025openhands}. Agents receive the complete pinned source repository and public
behavioral contract, while benchmark tests, reference implementations, and
benchmark-provided source-location hints are withheld.

Every run in the main comparison and paired ablation uses a 3,600-second timeout,
a 131,072-token context window with 8,192 tokens reserved for output, and the
same maximum of 150 interaction steps. In the main comparison, Pro and Flash use
token-based context condensation; the remaining configurations use the default
OpenHands context handling. Work on agent--computer interfaces and harness
effects motivates reporting these execution choices as part of the evaluated
system~\cite{yang2024sweagent,yao2026harnessbench}. We therefore treat each
model--harness configuration as the experimental unit rather than attributing
observed differences to the base model alone.

\subsection{Functional Evaluation and Artifact Metrics}

Only the collected submission enters the evaluator environment; the source
repository is unavailable at runtime, networking is disabled, and dependencies
are fixed. Functional success follows Equation~\ref{eq:functional}. The
\textbf{Build} gate checks that the submitted package can be loaded under the
frozen evaluator environment; \textbf{Primary} and \textbf{Extended} exercise
withheld behavioral obligations from the same public contract; and
\textbf{Isolation} checks prohibited source dependencies and protected-path
access.

\paragraph{Successful-artifact characteristics.}
For a functionally passing artifact $A$ and its frozen feasible reference
$A_{\mathrm{ref}}$, we define Reference-Relative Extraction Size as
\begin{equation}
  \operatorname{RRES}(A)=
  \frac{\operatorname{PythonLOC}(A)}
       {\operatorname{PythonLOC}(A_{\mathrm{ref}})}.
  \label{eq:rres}
\end{equation}
RRES characterizes artifact size relative to one feasible reference; it does
not measure minimality.

We additionally report \emph{Copy}, the fraction of submitted Python source
matched to the source repository in normalized runs of at least three non-empty,
non-comment lines. Copy measures syntactic overlap rather than semantic
equivalence, a distinction established in code-clone
research~\cite{roy2009clones}, and can miss rewritten equivalents.
Both RRES and Copy are reported only after functional success and
are not incorporated into the functional score.
\flb does not impose a minimal-extraction requirement: functional correctness
is determined by the destination behavior and source independence, while RRES
and Copy separately characterize the resulting implementation footprint.

\paragraph{Execution effort.}
Table~\ref{tab:main-results} reports recorded assistant-step counts and total
token usage for all assigned tasks. Total tokens are computed as input plus
output tokens for each run.

\subsection{Comparison Scope and Statistical Analysis}

The main comparison contains one retained outcome for each of six configurations
on each of the same 150 tasks, yielding 900 model--task outcomes. All assigned
tasks remain in the denominator, including runs that produce no usable
submission. Functional correctness is determined by the frozen evaluator rather
than by the agent runner's terminal status.

We treat the resulting pass rates as observed performance of the evaluated
model--harness configurations. Because each configuration contributes one
retained outcome per task, the main comparison does not estimate run-to-run
stochastic variance.

The source-evidence ablation compares Full Source and Contract Only on the same
tasks using paired pass-rate differences, paired bootstrap confidence intervals,
and exact conditional McNemar tests~\cite{fagerland2013mcnemar}, with Holm
correction across configurations~\cite{holm1979multiple}. Intervals use
100{,}000 paired task-bootstrap resamples and describe variation across task
pairs rather than repeated agent runs.

For successful-artifact comparisons, we retain tasks with successful artifacts
from at least two configurations and fit task fixed-effects models:
\begin{equation}
\begin{aligned}
  \log_2\!\left(\operatorname{RRES}_{t,m}\right)
    &= \alpha_t^{\mathrm{RRES}} + \beta_m^{\mathrm{RRES}}
       + \epsilon_{t,m}^{\mathrm{RRES}}, \\
  \operatorname{Copy}_{t,m}
    &= \alpha_t^{\mathrm{Copy}} + \beta_m^{\mathrm{Copy}}
       + \epsilon_{t,m}^{\mathrm{Copy}},
\end{aligned}
\label{eq:adjusted-footprint}
\end{equation}
with $\sum_m\beta_m^{\mathrm{RRES}}=\sum_m\beta_m^{\mathrm{Copy}}=0$.
Here $t$ indexes tasks and $m$ indexes configurations. Ordinary least squares
gives each successful artifact equal weight. We report the task-adjusted RRES
ratio $2^{\beta_m^{\mathrm{RRES}}}$ and the Copy difference
$100\beta_m^{\mathrm{Copy}}$ in percentage points. Their reference centers,
$1\times$ and 0 pp, represent the sum-to-zero center of configuration effects;
the RRES ratio is not an adjusted median or a raw size ratio to the reference
implementation.

Pointwise 95\% percentile intervals use 10{,}000 accepted task-cluster bootstrap
resamples~\cite{field2007cluster}, retaining all successful artifacts of each sampled task and refitting
and centering both models on every resample. We check that the six-configuration
co-success graph remains connected; a disconnected draw would be rejected and
redrawn. Sensitivity analyses give each artifact weight $1/k_t$, where $k_t$ is
the number of successful configurations on task $t$, or resample whole source
repositories while retaining artifact weights from the main analysis. Each uses
10{,}000 accepted resamples; no disconnected draws occurred in any scheme.
These intervals describe variation across task or repository clusters, not
repeated agent runs. The analysis is descriptive and conditional on success;
it does not estimate causal configuration effects or a global quality ranking.

\subsection{Paired Source-Evidence Ablation}
\label{sec:source-ablation-protocol}

The main benchmark intentionally provides the intact source repository because
feature lifting studies reconstruction from existing implementation evidence. To
test whether that evidence materially affects success, we conduct a paired
ablation that removes the repository while keeping the requested behavior fixed.

We select a fixed 40-task subset from 38 repositories, containing 15 Direct,
20 Adapted, and 5 Composite tasks. The subset was stratified by lift type and
frozen before inspecting ablation outcomes. GPT-5.6 Luna,
DeepSeek V4 Pro, and Qwen3.6-35B-A3B are each evaluated under two conditions
on the same tasks, yielding 240 retained outcomes. These are new paired runs
rather than reused outcomes from the main comparison.

In the \emph{Full Source} condition, the agent receives the complete pinned
repository together with the public behavioral contract. In \emph{Contract
Only}, the repository is omitted while the same contract and permitted
dependency information are retained; retrieving the original project from the
network, installed packages, or caches is prohibited. Benchmark tests, reference
implementations, and source-location hints remain hidden in both conditions. The
intervention therefore removes repository evidence as a whole, including source
code, upstream tests, documentation, and packaged resources, rather than
isolating source-code text alone.

Within each configuration, the paired conditions use the same model--harness
settings and source-free evaluator. The Pro configuration in this ablation does
not apply context compression. We compare functional success using
paired pass-rate differences, paired bootstrap confidence intervals, and exact
McNemar tests with Holm correction across the three configurations. Missing
submissions are counted as functional failures.

\subsection{Trace-Based Source-Exposure Measurement}
\label{sec:source-exposure-protocol}

To distinguish source availability from observable source use, we analyze the
900 saved trajectories from the main comparison. For
each task, we statically map its declared source entrypoints to files in the
pinned repository. This resolves 607 of 641 declared entrypoints; at least one
entrypoint is resolved for 145 of the 150 tasks.

We call an event a \emph{confirmed source read} when a successful tool read
returns content matching at least two adjacent normalized source lines from a
mapped file. Filename mentions, unsuccessful reads, and traceback text do not
qualify. Unresolved mappings remain unconfirmed rather than being treated as
absence of source inspection. This measure therefore captures observable
exposure to entrypoint-associated source-file content, not whether the agent
located the complete implementation, followed all supporting dependencies, or
understood the returned code. We use it as a diagnostic of whether behavioral
failures can persist despite confirmed reads of entrypoint-associated
source-file content.

\paragraph{Qualitative analysis of source-exposed Primary/Extended-first failures.}
We inspect a fixed, purposively diverse sample of 40 source-exposed
Primary/Extended-first failures, with one case per task. Selection uses fixed
configuration quotas and coverage of functional families, lift types, and prior
review status, with deterministic hash-based tie-breaking. The sample spans all
six configurations, ten functional families, and all three lift types.

Two authors independently reviewed the public contract, pinned upstream
implementation, submitted artifact, and evaluator outcome for each case, and
reconciled disagreements to derive a qualitative coding framework. We synthesize
the reviewed evidence into nonexclusive themes that relate observable obligations
to supporting software relationships and distinguish upstream preservation from
destination adaptation. Cases not captured by recurring themes are retained as other or case-specific
observations; the analysis identifies qualitative patterns rather than their
prevalence. The replication materials preserve the complete case-to-theme
mapping and review history. Three illustrative cases are paired with successful
artifacts under the same contracts.

\subsection{Retrospective Execution-Effort Analysis}
\label{sec:execution-effort-protocol}

We recover token usage from provider audit logs and persisted per-response
metrics, including agent and context-compression calls. Shared parent/child
records are deduplicated by response ID. Total token usage is computed as the sum of input and output tokens across all
model calls in each run.

For trajectory-level analysis, we evaluate reconstructed artifacts at saved
tool-completion boundaries and identify the first \emph{observed passing
checkpoint}. With cumulative tokens $T_r^{\mathrm{obs}}$ at this checkpoint and
run total $T_r$, the post-sufficiency fraction is
$\mathrm{PSF}_r=(T_r-T_r^{\mathrm{obs}})/T_r$.
Token timing requires one-to-one response-ID alignment with the event stream,
including compression responses, and agreement with audit counts and timing.
We separately count distinct subsequent responses recorded as actions or
messages in the main agent event stream, excluding the response that produced
the passing artifact. This count
excludes compression responses and unobserved HTTP retries. Both measures
condition on final success, evaluated checkpoint chains, and matching
reconstructed and re-evaluated final submissions, but use independent available
samples. Values are computed per run before aggregation. These retrospective
checkpoints do not establish a stopping signal available to the agent.

\section{Benchmark Results}
\label{sec:results}

We organize the core results around three questions about agent-mediated
software reuse: \textbf{RQ1} measures functional success overall and across reuse
structures; \textbf{RQ2} tests the effect of repository evidence through paired
ablation; and \textbf{RQ3} examines behavioral preservation failures despite
confirmed reads of entrypoint-associated source-file content. Secondary analyses
characterize execution after an observed passing checkpoint and the
implementation footprints of successful packages.

% BEGIN GENERATED TABLE: main
\begin{table}[htb]
  \caption{Main comparison on the common 150-task benchmark. Pass@1 and Steps
  use all assigned runs; RRES and Copy use only functionally successful artifacts.}
  \label{tab:main}

  \centering
  \footnotesize
  \setlength{\tabcolsep}{1.6pt}
  \renewcommand{\arraystretch}{1.12}

  \begin{tabular*}{\linewidth}{
    @{\extracolsep{\fill}}
    >{\raggedright\arraybackslash}p{2.65cm}
    r
    rrr
    rrr
    rr
    rr
    @{}
  }
    \toprule

    &
    \multicolumn{1}{c}{Correctness} &
    \multicolumn{3}{c}{RRES} &
    \multicolumn{3}{c}{Copy} &
    \multicolumn{2}{c}{Steps} &
    \multicolumn{2}{c}{Tokens ($10^4$)} \\

    \cmidrule(lr){2-2}
    \cmidrule(lr){3-5}
    \cmidrule(lr){6-8}
    \cmidrule(lr){9-10}
    \cmidrule(lr){11-12}

    Configuration
      & \shortstack{Pass@1\\$n$ (\%)}
      & Mean
      & Median
      & $[Q_1,Q_3]$
      & Mean
      & Median
      & $[Q_1,Q_3]$
      & Median
      & P90
      & Median
      & P90 \\

    \midrule
    \mbox{DeepSeek V4 Pro} & \textbf{115 (76.7)} & 1.78 & 0.99 & [0.77, 1.29] & 0.78 & 0.96 & [0.72, 0.99] & 41.5 & 80.1 & 8.80 & 22.53 \\
    \mbox{DeepSeek V4 Flash} & 108 (72.0) & 1.73 & 1.00 & [0.89, 1.22] & 0.81 & 0.97 & [0.84, 0.99] & 63.0 & 116.2 & 12.36 & 26.95 \\
    \mbox{GPT-5.6 Luna} & 102 (68.0) & 0.95 & 0.81 & [0.31, 1.00] & 0.41 & 0.16 & [0.05, 0.97] & 35.5 & 74.4 & 148.37 & 1134.94 \\
    \mbox{GLM-5.3-Flash} & 68 (45.3) & 2.47 & 1.01 & [0.92, 1.16] & 0.84 & 0.95 & [0.81, 0.98] & 85.5 & 127.0 & 701.85 & 1879.19 \\
    \mbox{Qwen3.6-35B-A3B} & 63 (42.0) & 1.56 & 0.97 & [0.65, 1.23] & 0.64 & 0.76 & [0.33, 0.97] & 41.5 & 106.6 & 181.86 & 552.10 \\
    \mbox{GPT-OSS 120B} & 36 (24.0) & 1.18 & 0.98 & [0.29, 1.65] & 0.31 & 0.18 & [0.01, 0.53] & 19.0 & 61.1 & 52.60 & 187.73 \\

    \bottomrule
  \end{tabular*}

\end{table}
% END GENERATED TABLE: main

\subsection{RQ1: How Does Feature-Lifting Success Vary Across Configurations and Reuse Structures?}
\label{sec:rq1}

Table~\ref{tab:main} reports the main comparison over the same 150 tasks for all
six configurations. DeepSeek V4 Pro achieves the highest observed functional
success, passing 115/150 tasks (76.7\%), followed by Flash at 108/150 (72.0\%)
and Luna at 102/150 (68.0\%). GLM, Qwen, and OSS pass 68 (45.3\%), 63 (42.0\%),
and 36 (24.0\%) tasks, respectively.

These rates establish the observed range of successful reconstruction under a
common reuse boundary and information policy. The structure breakdown below
examines how those outcomes vary with the transformation requested by each task.
Because each model--task cell contains one retained run, the values describe
observed performance of the evaluated model--harness configurations rather than
run-to-run success probabilities.




\paragraph{Observed variation across task structures.}
Table~\ref{tab:structure} reports pass rates within the three lift types
and four overlapping entanglement-mechanism groups. Across all six configurations, observed pass rates are lower on Adapted
and Composite tasks than on Direct tasks.
For these configurations, success declines when reuse requires changing or
composing the upstream capability boundary rather than directly preserving an
identifiable capability. For example, Pro passes 91.1\% of Direct tasks,
72.4\% of Adapted tasks, and 50.0\% of Composite tasks.
These descriptive comparisons do not isolate a causal difficulty effect:
category membership co-varies with other task properties, and the Composite
group contains only 18 tasks.
The mechanism groups also overlap substantially: a task may involve both shared
state and environment-dependent resources, for example. Their marginal pass rates
cannot identify the effect of an individual mechanism or explain which obligation
was violated in a particular failed artifact. RQ3 examines that distinction at
the level of concrete contracts and implementations.


% BEGIN RESULTS EVIDENCE: structure
\begin{table}[htb]
  \centering
  \footnotesize
  \caption{Observed functional pass rates (\%) by task structure. $n$ is the row denominator. Lift types partition the tasks; mechanism groups overlap.}
  \label{tab:structure}
  \setlength{\tabcolsep}{4.2pt}
  \renewcommand{\arraystretch}{1.10}
  \begin{tabular}{@{}lrrrrrrr@{}}
    \toprule
    Category & $n$ & Pro & Flash & Luna & GLM & Qwen & OSS \\
    \midrule
    Direct & 56 & 91.1 & 89.3 & 80.4 & 71.4 & 55.4 & 28.6 \\
    Adapted & 76 & 72.4 & 67.1 & 65.8 & 31.6 & 35.5 & 21.1 \\
    Composite & 18 & 50.0 & 38.9 & 38.9 & 22.2 & 27.8 & 22.2 \\
    \midrule
    Code dependencies & 139 & 79.9 & 73.4 & 69.1 & 46.0 & 42.4 & 22.3 \\
    Data and state & 127 & 80.3 & 75.6 & 71.7 & 48.0 & 44.1 & 23.6 \\
    Framework mechanisms & 71 & 78.9 & 73.2 & 66.2 & 40.8 & 43.7 & 16.9 \\
    Environment and resources & 49 & 75.5 & 71.4 & 63.3 & 32.7 & 34.7 & 24.5 \\
    \bottomrule
  \end{tabular}
\end{table}
% END RESULTS EVIDENCE: structure

\subsection{RQ2: Does Repository Evidence Improve Behavioral Reconstruction?}
\label{sec:rq2}
\label{sec:source-ablation-results}

Figure~\ref{fig:source-evidence} visualizes the paired effects, while
Table~\ref{tab:paired-ablation} reports the exact paired statistics.
Repository evidence improves Luna's and Qwen's functional success, with both
paired differences remaining significant after Holm correction. Pro also shows
a positive difference, but its Contract Only condition includes 18 runs without
usable submissions. Pro's paired effect therefore reflects both behavioral
reconstruction and successful artifact delivery.

\begin{figure}[htb]
  \centering
  \begin{minipage}[c]{0.48\linewidth}
    \centering
    \includegraphics[width=\linewidth]{figures/fig5a_pass_rate.pdf}
  \end{minipage}
  \hfill
  \begin{minipage}[c]{0.48\linewidth}
    \centering
    \includegraphics[width=\linewidth]{figures/fig5b_paired_gain.pdf}
  \end{minipage}
  \caption{Source-evidence ablation on the same 40 tasks per configuration:
  (a) functional pass rates under Full Source and Contract Only;
  (b) Full Source minus Contract Only effects with 95\% paired task-bootstrap
  intervals. The 18 Pro Contract Only runs without submissions are counted
  as failures.}
  \Description{Two panels compare Full Source and Contract Only for GPT-5.6
  Luna, DeepSeek V4 Pro, and Qwen. Luna and Qwen have substantially higher pass
  rates with repository evidence. Pro also has a large difference, with
  18 runs without usable submissions counted as Contract-Only failures.}
  \label{fig:source-evidence}
\end{figure}

% BEGIN RESULTS EVIDENCE: paired-ablation
\begin{table}[htbp]
  \centering
  \small
  \caption{Paired source-evidence outcomes and statistics on 40 tasks per configuration. Full-only and Contract-only count discordant pairs. $\Delta$ is Full Source minus Contract Only (pp); intervals are 95\% paired task-bootstrap intervals, and $p_{\mathrm{Holm}}$ gives exact McNemar p-values adjusted across the three configurations.}
  \label{tab:paired-ablation}
  \setlength{\tabcolsep}{3.6pt}
  \renewcommand{\arraystretch}{1.10}
  \begin{tabularx}{\linewidth}{@{}>{\raggedright\arraybackslash}X>{\centering\arraybackslash}p{0.095\linewidth}>{\centering\arraybackslash}p{0.095\linewidth}>{\centering\arraybackslash}p{0.07\linewidth}>{\centering\arraybackslash}p{0.095\linewidth}>{\centering\arraybackslash}p{0.05\linewidth}>{\centering\arraybackslash}p{0.145\linewidth}>{\centering\arraybackslash}p{0.13\linewidth}@{}}
    \toprule
    & \multicolumn{2}{c}{\textbf{Passes}} & \multicolumn{2}{c}{\textbf{Discordant pairs}} & \multicolumn{3}{c}{\textbf{Paired effect}} \\
    \cmidrule(lr){2-3}\cmidrule(lr){4-5}\cmidrule(l){6-8}
    \textbf{Configuration} & \shortstack{\textbf{Full}\\\textbf{source}} & \shortstack{\textbf{Contract}\\\textbf{only}} & \shortstack{\textbf{Full}\\\textbf{only}} & \shortstack{\textbf{Contract}\\\textbf{only}} & \shortstack{\boldmath$\Delta$\\\textbf{(pp)}} & \shortstack{\textbf{95\% CI}\\\textbf{(pp)}} & \boldmath$p_{\mathrm{Holm}}$ \\
    \midrule
    GPT-5.6 Luna & 23/40 & 9/40 & 17 & 3 & 35.0 & [15.0, 55.0] & $5.15\times10^{-3}$ \\
    DeepSeek V4 Pro & 25/40 & 6/40 & 20 & 1 & 47.5 & [30.0, 65.0] & $6.29\times10^{-5}$ \\
    Qwen3.6-35B-A3B & 12/40 & 1/40 & 12 & 1 & 27.5 & [12.5, 42.5] & $5.15\times10^{-3}$ \\
    \bottomrule
  \end{tabularx}
\end{table}
% END RESULTS EVIDENCE: paired-ablation

\paragraph{The gain reflects behavioral reconstruction, not merely delivery.}
For both Luna and Qwen, every task that succeeds only under Full Source has a
Contract-Only counterpart that still produces an artifact but fails a behavioral
gate. Repository evidence
therefore helps agents satisfy behavioral obligations that are missed even when
Contract Only already yields a deliverable package.

The benefit is substantial but incomplete. With Full Source available, Luna and
Qwen still exhibit Primary/Extended-first failures, showing that access to an existing
implementation does not guarantee successful reconstruction across the new
package boundary. Because the intervention removes the repository as a whole,
these results establish the value of repository evidence---including source
code, upstream tests, documentation, and resources---rather than the isolated
causal contribution of source-code text.
The source-exposure analysis in RQ3 complements this intervention. Together,
the analyses support
treating repository material as useful reconstruction evidence whose application
still requires checking the destination contract. Neither analysis establishes
that reading an entrypoint completes the recovery of all required behavior.

\subsection{RQ3: How Can Behavioral Preservation Fail Despite Confirmed Reads of Entrypoint-Associated Source-File Content?}
\label{sec:rq3}

We first locate the observed boundary at which each unsuccessful run ceases to
satisfy the evaluation protocol. We classify each run as Pass, no usable
submission, or by its first failed evaluator gate in the sequence
Build--Primary--Extended--Isolation. This classification identifies where
failure becomes observable; it does not by itself establish the underlying
cause.

Figure~\ref{fig:failures} shows a clear pattern among the strongest
configurations. Pro and Flash produce usable submissions for all 150 tasks and
have no Build-first failures. Of their 77 combined unsuccessful runs, 76 first
fail at Primary or Extended, and only one first fails at Isolation. Thus, once
these configurations deliver a package, the dominant remaining boundary is
Primary/Extended evaluation rather than package production or loading.

Weaker configurations exhibit a broader mixture of failure modes. GLM and Qwen
produce no usable submission for 38 and 25 tasks, respectively, while OSS
usually delivers an artifact but has 18 Build-first and 93
Primary/Extended-first failures. Low overall success therefore need not reflect
the same operational bottleneck across configurations.

\begin{figure}[htbp]
  \centering
  \includegraphics[width=\linewidth]{figures/fig4.pdf}
  \caption{Observed outcomes across 150 tasks per configuration, classified as
  Pass, no usable submission, or by the first failed evaluation gate in
  Build--Primary--Extended--Isolation order.}
  \Description{Horizontal stacked bars show Pass, no usable submission,
  Build-first, Primary-first, Extended-first, and Isolation-first outcomes for
  six configurations over 150 tasks each.}
  \label{fig:failures}
\end{figure}

\paragraph{Primary/Extended-first failures often persist after source exposure.}
A Primary/Extended-first failure could arise simply because the agent never inspected
implementation evidence associated with the requested capability. The saved
trajectories show that this explanation is insufficient for many runs. Among the
303 Primary/Extended-first failures across all six configurations, 241 (79.5\%)
contain confirmed reads of entrypoint-associated source-file content
(Table~\ref{tab:source-exposure}). The median first confirmed read among these runs occurs at tool action
five. For Pro and Flash alone, 63 of their 76 Primary/Extended-first failures contain
such exposure.

These observations show that opening entrypoint-associated source files is not
sufficient for successful feature lifting. A confirmed read does not establish
that the agent identified all supporting dependencies or understood all relevant
behavioral obligations, so the result should not be interpreted as a causal diagnosis
of why individual runs fail. It nevertheless rules out a simple account in
which most Primary/Extended-first failures occur before any associated source content is
observed.

% BEGIN VERIFIED SOURCE EXPOSURE TABLE
\begin{table}[htbp]
  \centering
    \caption{Confirmed reads of entrypoint-associated source-file content across
    900 main-comparison outcomes. Read fractions use all runs in each row;
    first-read action-step medians use only confirmed reads.}
    \label{tab:source-exposure}

    \centering
    \footnotesize

    \begin{tabular}{@{}lrrr@{}}
      \toprule
      Outcome
        & Runs
        & \shortstack{Confirmed read\\$n/N$ (\%)}
        & \shortstack{Median first\\read step} \\
      \midrule
      Pass & 492 & 363/492 (73.8) & 5 \\
      Primary/Extended-first & 303 & \textbf{241/303 (79.5)} & 5 \\
      Delivery/build & 101 & 51/101 (50.5) & 5 \\
      Isolation-first & 4 & 4/4 (100.0) & 10.5 \\

      \bottomrule
    \end{tabular}

\end{table}
% END VERIFIED SOURCE EXPOSURE TABLE

\paragraph{Recurring themes across the reviewed sample.}
Across the 40 reviewed cases, we identified three recurring qualitative
patterns:
\emph{changed ambient assumptions}, \emph{incompatible supporting representations
and pipelines}, and \emph{destination-specific obligations requiring adaptation}.
Platform-dependent path helpers in \texttt{platformdirs} and
\texttt{jupyter\_core}, and a construction-time scheduling bound in
\texttt{apscheduler}, illustrate the first theme. Grouped diagnostics in
\texttt{cerberus} and normalization before hashing in \texttt{requests\_cache}
illustrate the second. Resource containment, lock-file cleanup, short-circuit
execution, and configuration parsing illustrate the third.

The themes do not exhaust the sample: other cases involve local behavioral
differences, missing members, an artifact constraint, or insufficient evidence
for one of these relationships. We retain them as other or case-specific
observations rather than forcing a theme assignment. The mapping in the replication materials
records every case and reports counts only within this purposive sample.
Figure~\ref{fig:failure-analysis} illustrates one instance of each recurring
theme with a paired failed/successful implementation. Each pair shares the same
archived public contract and pinned source; the successful counterpart passes
all four functional gates.

\begin{figure}[htbp]
  \centering
  \includegraphics[width=\linewidth]{figures/fig6.png}
  \caption{Illustrative supporting relationships in Primary/Extended-first
failures with confirmed source-file content reads. Each row connects upstream
implementation evidence to a destination obligation and contrasts failed and
passing agent artifacts under the same public contract. Code and structures are
schematic; the paired artifacts may differ in multiple implementation choices.}
  \Description{Three paired cases illustrate changed ambient assumptions in
  platformdirs, incompatible supporting representations in cerberus, and
  destination-specific adaptation in importlib resources. A dashed line marks
  the new package boundary between upstream evidence and destination obligations.
  Red and green panels contrast failed and passing reconstructions and their
  evaluated outcomes.}
  \label{fig:failure-analysis}
\end{figure}

\paragraph{An explicit input can invalidate an ambient assumption.}
In \texttt{platformdirs}, the destination contract turns platform selection into
an explicit input. The failed artifact selects a Windows branch but still uses
the host's path-joining operations. The successful counterpart uses a
Windows-aware helper for those paths. Selecting the requested branch is therefore
not enough: the helpers below that branch must honor the same platform choice.
This case illustrates adaptation of an upstream assumption rather than simply
omitting an upstream API.

\paragraph{A preserved producer needs a compatible consumer.}
For \texttt{cerberus}, nested validation and diagnostic rendering share a grouped
error representation. The failed artifact produces child and group records,
but its formatter lacks compatible group-aware handling. The successful artifact
connects group production to recursive rendering. The public obligation concerns
structured nested diagnostics, so retaining validation logic alone does not
ensure that the observable result is preserved. Here the relationship supports
an inherited upstream behavior.

\paragraph{A destination restriction may require an additional adapter.}
For \texttt{importlib\_resources}, the destination contract requires traversal to
remain within an exposed resource tree. The failed artifact delegates access to
an unguarded resource object; the successful counterpart routes access through
a wrapper that resolves child names within the exposed tree. Containment is a
destination requirement, not an upstream guarantee lost during copying.
Passing the recorded traversal checks does not establish general filesystem
security.

Together, these comparisons illustrate how source-backed components can be
present while their assumptions or interactions remain incompatible with
destination behavior. They motivate checking the relationships that realize a
public obligation and identifying where adaptation changes the assumptions
under which source code is reused.

\FloatBarrier

\subsection{Secondary Analyses: Execution and Implementation Footprints}
\label{sec:secondary-analyses}
These analyses characterize execution and the packages that satisfy the
contract; they complement the core findings on repository evidence and
behavioral reconstruction.

\subsubsection{Execution after an observed passing checkpoint.}
\label{sec:execution-effort-results}

Final pass rates do not reveal when a sufficient artifact first appears.
We examine the separation between the first evaluator-observed sufficient
artifact and run termination. Observed behavioral sufficiency means passing
the frozen evaluator at a reconstructed checkpoint; it establishes neither
complete equivalence to the upstream feature nor a completion signal available
to the agent. The replication materials report outcome-level token distributions
for 764 recoverable ledgers; their incomplete coverage and mixed directions
across configurations do not support an efficiency ranking.

\begin{figure}[ht]
  \centering
  \begin{minipage}[c]{0.48\linewidth}
    \centering
    \includegraphics[width=\linewidth]{figures/fig4a_token_fraction.pdf}
  \end{minipage}
  \hfill
  \begin{minipage}[c]{0.48\linewidth}
    \centering
    \includegraphics[width=\linewidth]{figures/fig4b_responses.pdf}
  \end{minipage}
  \caption{Execution after the first reconstructed passing checkpoint among
  recoverable, finally successful runs: (a) fraction of input-plus-output tokens,
  including cached input (317 runs); (b) subsequent main-agent responses
  (403 runs). Panels use separate available samples, with configuration counts
  shown below the axes. Boxes show medians and quartiles, whiskers extend to
  1.5 IQR, and points denote runs.}
  \Description{Two panels show distributions for six configurations. The token panel contains 98, 93, 43, 27, 25, and 31 runs; the response panel contains 102, 96, 75, 46, 50, and 34 runs, respectively.}
  \label{fig:execution-effort}
\end{figure}

\paragraph{Execution continues after an observed passing artifact.}
Figure~\ref{fig:execution-effort} shows continued execution across all six
configurations. The token panel contains 317 successful runs, while the response
panel contains 403; each panel reports its own sample sizes. Median
post-checkpoint token fractions are 66.9\% for Pro, 73.0\% for Flash, 58.1\%
for Luna, 67.1\% for GLM, 53.7\% for Qwen, and 37.1\% for OSS. The corresponding
configuration medians in the response panel are 20.5, 32.5, 11, 39.5, 17.5, and
5 additional responses. Substantial continuation therefore appears in both
token consumption and agent interaction, rather than only in final run totals.
The recoverable task subsets differ across configurations and do not establish
a task-adjusted efficiency ranking.



On the same token samples, using the first checkpoint after which all remaining
reconstructed versions pass gives medians of 66.6\%, 71.1\%, 57.9\%, 54.8\%,
52.1\%, and 37.1\%, respectively. Continuation thus persists when transient
passes followed by regressions are excluded. It need not be unproductive: in
an inspected Pro JSONPath trajectory, subsequent actions run tests, inspect
dependencies, and clean generated files without changing the evaluated artifact.
Repeated context consumption also matters: under the uncached-input accounting
of Table~\ref{tab:main}, Pro and Flash have median post-checkpoint fractions of
34.2\% and 43.9\% on their respective token samples. The observed continuation
motivates studying verification and stopping decisions; it does not directly
quantify avoidable tokens, cost, or computation.

\subsubsection{Implementation footprints of successful artifacts.}
\label{sec:footprint-results}

Successful artifacts can satisfy the same contract with different implementation
sizes and source overlap. To account for differing task sets, the adjusted
analysis includes 485 successful artifacts from 115 tasks with at least two
successful configurations, spanning 97 repositories. Seven singleton successes
are excluded. Equation~\ref{eq:adjusted-footprint} defines the adjustment;
Figure~\ref{fig:matched-footprint} uses a symmetric configuration-effect center
and pointwise 95\% task-bootstrap intervals.

\begin{figure}[ht]
  \centering
  \begin{minipage}[c]{0.48\linewidth}
    \centering
    \includegraphics[width=\linewidth]{figures/fig8a_rres.pdf}
  \end{minipage}
  \hfill
  \begin{minipage}[c]{0.48\linewidth}
    \centering
    \includegraphics[width=\linewidth]{figures/fig8b_copy.pdf}
  \end{minipage}
  \caption{Task-adjusted footprints of 485 successful artifacts across
  115 tasks: (a) RRES ratios relative to the configuration-effect center;
  (b) centered Copy differences in percentage points. Dashed lines mark the
  centers (1$\times$ and 0 pp); intervals are pointwise 95\% task-cluster
  bootstrap intervals.}
  \Description{Two vertical bar charts compare Pro, Flash, Luna, GLM, Qwen,
  and OSS. The left panel shows task-adjusted implementation-size ratios with
  a reference line at one; the right shows source-overlap differences with a
  reference line at zero percentage points. Luna and OSS are below both
  centers; Qwen's source-overlap interval crosses zero. Each bar has a 95 percent
  task-bootstrap interval.}
  \label{fig:matched-footprint}
\end{figure}

Luna and OSS are below this center on both RRES and Copy, while Pro, Flash,
and GLM are above it. Qwen has an RRES ratio above one, but its Copy interval
spans zero. These are footprint profiles, not a complete pairwise ranking.
The replication materials provide exact estimates and sample counts; the accompanying
weighting and repository-cluster checks preserve the qualitative profiles.
The near-equal RRES estimates for Luna and OSS do not establish their ordering.



RRES and Copy measure size relative to a feasible reference and detected textual
source overlap. Neither measures minimality, maintainability, or semantic
superiority, and neither identifies a deliberate reuse strategy without process
evidence. These results characterize successful artifacts only and do not
extrapolate to tasks on which a configuration fails.

\FloatBarrier
\section{Discussion}
\label{sec:diagnostics}

\paragraph{Define reusable features by destination behavior.}
Feature lifting suggests that a reusable capability is better defined by its
observable obligations than by any particular source entrypoint or code slice.
Repository evidence can substantially aid reconstruction, yet locating and
reusing relevant implementation components does not by itself ensure that the
resulting package satisfies the destination contract. For extraction and
migration tasks, the behavioral contract therefore provides a more appropriate
reuse boundary than source structure alone.

\paragraph{Verify the relationships that support observable behavior.}
Behavioral obligations may depend on helpers, shared representations, state,
resources, or other relationships beyond the apparent entrypoint. Feature
extraction and migration tools should therefore reason not only about code
dependencies, but also about the relationships needed to realize externally
observable behavior at the new boundary. The objective is behavioral
correctness rather than structural fidelity to the donor implementation.

\paragraph{Make adaptation assumptions explicit across software boundaries.}
Source code expresses behavior under assumptions made by its original
environment, while a destination may change interfaces, platform conditions,
resource policies, or other constraints. Such changes should be treated
explicitly as adaptation obligations rather than assumed to follow automatically
from upstream behavior. Making changed assumptions visible can help developers
and tools distinguish inherited behavior from destination-specific requirements
and verify each accordingly.

\paragraph{Secondary implications for agent execution and evaluation.}
Passing artifacts can appear before run termination, motivating better
verification and stopping mechanisms, although continued execution is not
necessarily avoidable because agents do not observe the offline evaluator.
Evaluation should also distinguish package delivery, behavioral correctness,
source independence, and implementation footprint; artifact size and source
overlap characterize successful implementations but do not substitute for
behavioral verification.

\section{Threats to Validity}
\label{sec:threats}

\paragraph{Benchmark scale.}
The current version of \flb contains 150 tasks from 126 repositories, which
provides diverse coverage but cannot exhaust the space of real-world feature-lifting
scenarios. Future versions will expand the benchmark with more tasks, repositories,
and feature types to improve coverage and support broader evaluation.

\paragraph{Language coverage.}
The current benchmark focuses on Python. Feature lifting may exhibit different
characteristics in languages with different type systems, build systems, and
dependency mechanisms. We plan to extend \flb to additional ecosystems,
including Java and Go, to study how well the findings generalize across
programming languages.

\section{Related Work}
\label{sec:related-work}

Feature lifting connects repository-level coding evaluation with software reuse
and behavior-preserving transformation. We organize the comparison around the
implementation evidence available to an agent and the boundary at which its
deliverable must remain correct.

\subsection{Repository-Level Coding and Feature Development}

SWE-bench evaluates repository patches that resolve real GitHub
issues~\cite{jimenez2024swebench}, and SWE-Bench Pro extends this direction to
longer-horizon software engineering tasks~\cite{deng2025swebenchpro}.
Feature-development benchmarks broaden the objective beyond issue resolution.
FeatBench evaluates implementation from natural-language requirements without
code hints, including whether changes preserve existing
features~\cite{chen2026featbench}. FeatureBench derives tasks through test-driven
dependency tracing and evaluates both incremental development in an incomplete
repository and construction of the requested functionality from
scratch~\cite{zhou2026featurebench}.

Recent benchmarks extend evaluation toward standalone repository construction.
NL2Repo-Bench asks agents to construct an installable Python library from a
natural-language requirements document and an empty
workspace~\cite{ding2026nl2repo}. RepoZero formulates repository generation as
black-box reproduction from API specifications without access to the original
library implementation~\cite{zhang2026repozero}. These settings study construction
from specifications when the donor implementation is unavailable.

Together, these benchmarks cover issue resolution, feature development, and
repository construction under different information boundaries. \flb occupies
a distinct source-backed reuse setting: an intact donor repository is available
as implementation evidence during construction, but a bounded capability must
remain correct after crossing into an independent package, evaluated without
runtime access to the donor.

Table~\ref{tab:task-comparison} summarizes these task interfaces. The comparison
concerns what each setting asks an agent to do, not relative difficulty.

\begin{table}[htbp]
  \caption{Task-interface comparison with closely related construction, reuse,
  and refactoring settings.}
  \label{tab:task-comparison}
  \centering
  \small
  \setlength{\tabcolsep}{4.2pt}
  \renewcommand{\arraystretch}{1.12}
  \begin{tabularx}{\linewidth}{
    @{}
    >{\raggedright\arraybackslash}p{0.23\linewidth}
    >{\centering\arraybackslash}X
    >{\centering\arraybackslash}p{0.22\linewidth}
    >{\centering\arraybackslash}p{0.22\linewidth}
    @{}
  }
    \toprule
    \textbf{Setting} & \textbf{Starting point} & \textbf{Deliverable} & \textbf{Evaluation boundary} \\
    \midrule
    SWE-bench / Pro
      & Repository + issue
      & Project patch
      & Revised project \\
    FeatBench
      & Base repository + requirements
      & Feature implementation
      & Feature + regressions \\
    FeatureBench L1
      & Incomplete repository + interface
      & Incremental implementation
      & Feature + regressions \\
    FeatureBench L2
      & Task/interface; no repository
      & Library from scratch
      & Generated library \\
    RepoZero
      & API specifications; no donor implementation
      & Repository from scratch
      & Black-box behavioral equivalence \\
    Donor--host transplantation
      & Donor + host + tests
      & Functionality in host
      & Integrated functionality \\
    SWE-Refactor
      & Target method + refactoring type + repository context
      & Refactored code
      & Project build, tests, and transformation checks \\
    \midrule
    \textbf{FeatureLiftBench}
      & \textbf{Intact source + behavioral contract}
      & \textbf{Independent package}
      & \textbf{Source-free functionality} \\
    \bottomrule
  \end{tabularx}
\end{table}

\subsection{Feature Location, Extraction, and Software Reuse}

Program slicing uses dependence analysis to retain statements relevant to a
specified criterion~\cite{weiser1984slicing}, while feature location identifies
source entities associated with a requested concern~\cite{dit2013featurelocation}.
Research on multi-dimensional separation of concerns addresses functionality
that cuts across module boundaries~\cite{tarr1999degrees}. These foundations
explain why recovering a reusable capability can require reasoning beyond its
entrypoint: the implementing statements and their supporting dependencies need
not form an existing module.

Software transplantation turns this recovery problem into executable reuse.
Barr et~al. combine analysis, testing, and search to extract donor functionality
and adapt it to a supplied host~\cite{barr2015transplantation}. Foundry extends
transplantation to software product lines; its prodScalpel tool uses tests and
annotations identifying features and implantation points to extract and
integrate functionality across codebases~\cite{souza2025foundry}.
\flb builds on this reuse objective through a benchmark for general-purpose
coding agents. It supplies the intact donor and a behavioral contract, while
leaving source localization and reconstruction to the agent. The destination is
an independent package, and the contract permits extraction, adaptation, or
compatible reimplementation. The contribution is an evaluation setting for this
end-to-end capability, complementary to techniques that automate particular
extraction and integration steps.

\subsection{Behavior-Preserving Refactoring and Transformation}

Classic refactoring work defines restructuring operations with preconditions
under which program behavior is preserved~\cite{opdyke1992refactoring}.
EM-Assist combines LLM suggestions with static analysis to
select Extract Method candidates and uses IDE machinery to perform the
transformation~\cite{pomian2024emassist}. SWE-Refactor evaluates atomic and
compound refactorings mined from real Java projects. Its verification combines
project compilation and tests with checks that the requested structural
transformation occurred~\cite{xu2026swerefactor}.

More recently, SWE Refactor Bench studies whole-repository stack migrations and
separates migration completeness from behavioral correctness: its evaluation
verifies that the requested migration occurred in addition to testing preserved
behavior~\cite{hong2026swerefactorbench}.

Feature lifting imposes no prescribed refactoring or stack transformation. Its
preservation obligation instead lies at the exported interface of a new package:
the agent must reconstruct enough supporting implementation for the specified
behavior to survive without runtime access to the donor repository. Structural
correspondence or minimal extraction is therefore not an acceptance requirement;
artifact size and source overlap characterize successful implementations rather
than define correctness.

\subsection{Repository Context and Evidence Utilization}

RepoCoder uses iterative retrieval and generation for repository-level code
completion~\cite{zhang2023repocoder}. RepoBench evaluates retrieval, completion,
and their combined pipeline~\cite{liu2024repobench}, while CrossCodeEval targets
completion requiring cross-file context~\cite{ding2023crosscodeeval}. These works
establish repository information as a resource for generation.
BeyondSWE extends evaluation to tasks requiring evidence beyond the target
repository. Its SearchSWE experiments report limited and uneven gains from search
augmentation and illustrate failures to integrate retrieved information with
local implementation and version constraints~\cite{chen2026beyondswe}.
\flb studies a complementary setting in which the relevant external evidence is
the intact donor repository itself. The source-evidence ablation measures its
contribution to reconstruction, and the failure analysis examines unmet
behavioral obligations after confirmed source reads, connecting evidence
utilization to the correctness of a source-independent artifact.

\section{Conclusion}

Repository evidence improves reconstruction, but source exposure does not ensure
that behavior survives a new package boundary. 79.5\% of Primary/Extended-first
failures contain confirmed reads of entrypoint-associated source-file content. Qualitative analysis of a purposively diverse 40-case sample identifies
recurring but limited patterns involving changed ambient assumptions, incompatible supporting
representations and pipelines, and destination-specific obligations requiring
adaptation. Paired implementations illustrate these patterns through platform,
diagnostic, and resource relationships. Successful lifting requires translating
source evidence into destination-contract behavior, including adaptations to
changed assumptions.
The benchmark provides an executable setting for studying that translation and
its verification across software boundaries.

\section*{Data Availability}
% TODO before submission: replace this statement with the exact anonymous
% artifact URL/identifier and confirm the listed contents.
The benchmark specifications, evaluation harness, analysis scripts, and
run-level artifacts needed to reproduce the reported results are included in the
anonymized replication package submitted with the paper.

\FloatBarrier
\bibliographystyle{ACM-Reference-Format}
\bibliography{references}

\end{document}
